% \section{Problem of Digital Conservation of Musical Instruments}
% \section{Open Science}
% \subsection{F.A.I.R.}
% \subsection{Open Science for Teaching}
% \subsection{Sustainable Software Engineering}
% \section{Digital Conservation}
% \subsection{Musical Instruments and Digital Cultural Heritage}
% \section{Thesis Structure}


% \section{Curated Decay}

% \begin{quote}
% As the chimney becomes less legible as an object of industrial heritage, it
% becomes possible to read other narratives out of its remains, to trace the
% granite blocks from their source in the Cornubian batholith to their temporary
% enrollment in this structure, to follow the ivy roots into the seams of the
% stone to learn how they find nourishment in mineral mortar, to envision a future
% in which the chimney no longer stands but something of its substance and its
% story persists nonetheless—to understand change not as loss but as a release
% into other states, unpre- dictable and open.
% \end{quote}

% \begin{quote}
% The prevailing preservation paradigm, which de- clares that certain objects must
% be retained for the benefit of future generations and asserts the moral
% imperative of material conservation, only emerged in the late nineteenth
% century.3 
% \end{quote}

% \begin{quote}
% Graham Fairclough writes, “The obsession with physical conservation became so
% embedded in twentieth century mentalities that it is no longer easy to separate
% an attempt to understand the past and its meaning from agonising about which
% bits of it to protect and keep. . . . The remains of the past . . . seem to
% exist only to be preserved.”5
% \end{quote}

% \begin{quote}
% Cornelius Holtorf asserts that processes of change and crea- tive transformation
% may actually help maintain a connection to the past rather than sever it.12
% \end{quote}